\section{Related Work}
\label{related}

\textbf{LLM Agents in Economic Settings.}
Prior work deploys LLMs in economic simulation---EconAgent \citep{li2024econagent} at the macroeconomic level, QuantAgent \citep{wang2024quantagent} and FinAgent \citep{zhang2024finagent} for single-agent trading, and the LLM Economist \citep{karten2025llm} for mechanism design in tax policy.
Vending-Bench \citep{backlund2025vendingbenchbenchmarklongtermcoherence} benchmarks single-agent business coherence over long horizons; Vending-Bench Arena \citep{vendingbencharena2025} extends this to competitive multi-agent play, finding that frontier models independently develop monopoly exploitation and de-facto price cartels.
We focus instead on the systemic failure modes that emerge in high-frequency multi-agent marketplaces---destructive price spirals and coordinated Sybil fraud---and on training agents to prevent them.

\textbf{Market Instability and Sybil Attacks.}
The 2010 Flash Crash showed that individually rational algorithmic agents can collectively trigger market-wide collapse via positive feedback loops~\citep{kirilenko2017flash,farmer2009economy,johnson2013abrupt}.
Calvano et al.~\citep{calvano2020artificial} found that Q-learning agents tacitly collude above the competitive equilibrium without explicit communication; we show LLM agents exhibit the \textit{opposite} pathology---destructive undercutting below unit cost.
On the fraud side, the Sybil attack~\citep{douceur2002sybil} and its marketplace variants (fake-review campaigns, identity cycling~\citep{luca2016fake,mayzlin2014promotional}) are well-studied; Dellarocas~\citep{dellarocas2006strategic} showed reputation manipulation is rational when identity cost is low---a condition LLMs satisfy trivially.
We are the first to study Sybil attacks executed by a single LLM coordinating semantically diverse but fraudulently equivalent listings across multiple identities.

\textbf{Multi-Agent Frameworks and Alignment.}
AutoGen~\citep{wu2023autogen}, AgentBench~\citep{liu2023agentbench}, WebArena~\citep{zhou2023webarena}, and MetaGPT~\citep{hong2023metagpt} study cooperative or task-completion agent behavior, not adversarial equilibrium dynamics.
Constitutional AI~\citep{bai2022constitutional} and RLHF~\citep{ouyang2022training} optimize per-interaction helpfulness, which does not capture \textit{systemic} economic safety: an agent offering the lowest price can be individually helpful yet collectively catastrophic.
Our approach uses LoRA-based SFT~\citep{hu2022lora} on market traces filtered by Economic Alignment Score, directly training agents to internalize market externalities rather than redesigning incentives.
